\chapter{Ideas and Endogenous Growth}\label{ch:b2:endogenous-growth}

In the Solow and \omterm{def:b2:ramsey-model:ramsey}{Ramsey models}, growth of income per person in the long run comes from technology, and technology is
left unexplained. This chapter models where it comes from: ideas, produced by people who work on them. The chart sets two
facts side by side: how fast American productivity has grown, decade by decade, and how many people the large economies
now employ in research.

\begin{openingchart}{Left: average annual growth of total factor productivity, United States, by decade, 1954--2023, per
  cent. Right: researchers in research and development per million people, 1996 to the latest year.}{Penn World Table
  (rtfpna); World Bank, World Development Indicators (SP.POP.SCIE.RD.P6); our computation}
\begin{tikzpicture}
\begin{axis}[omrecord, name=tfp, ybar, bar width=9pt, width=0.47\linewidth, height=0.44\linewidth, xmin=0.4,
    xmax=7.6, ymin=-0.5, ymax=1.6, xtick={1,...,7}, xticklabels={1954--63,1963--73,1973--83,1983--93,1993--03,
    2003--13,2013--23}, xticklabel style={rotate=60, anchor=east, font=\scriptsize}, ylabel={Per cent a year},
    xmajorgrids=false]
  \addplot[fill=omDef!50, draw=omDef] table[x=pos, y=growth] {figdata/bachelor-2/24-endogenous-growth/out/tfp.csv};
\end{axis}
\begin{axis}[omrecord, at={(tfp.east)}, anchor=west, xshift=2.3cm, width=0.44\linewidth, height=0.44\linewidth,
    xmin=1996, xmax=2023, ymin=0, ymax=10000, xtick={2000,2010,2020}, xticklabel style={/pgf/number format/1000 sep=},
    scaled y ticks=false, yticklabel style={/pgf/number format/1000 sep={\,}}, ylabel={Per million people},
    unbounded coords=jump, legend style={at={(0.5,-0.32)}, anchor=north, legend columns=3, font=\scriptsize}]
  \addplot[omDef, thick] table[x=year, y=USA] {figdata/bachelor-2/24-endogenous-growth/out/researchers.csv};
  \addplot[omThm, thick, dashed] table[x=year, y=DEU] {figdata/bachelor-2/24-endogenous-growth/out/researchers.csv};
  \addplot[omMeth, thick, dotted] table[x=year, y=JPN] {figdata/bachelor-2/24-endogenous-growth/out/researchers.csv};
  \addplot[black!70, thick, dashdotted] table[x=year, y=KOR] {figdata/bachelor-2/24-endogenous-growth/out/researchers.csv};
  \addplot[omProp, thick] table[x=year, y=CHN] {figdata/bachelor-2/24-endogenous-growth/out/researchers.csv};
  \legend{United States, Germany, Japan, Korea, China}
\end{axis}
\end{tikzpicture}
\end{openingchart}

% figdata: tfp (1.33, 0.86, -0.23, 0.75, 0.96, 0.38, 0.53); res_summary (USA 3207 -> 4937 in 2022, ratio 1.54; DEU 2.11; JPN 1.14; KOR 4.33; CHN 4.71)
American productivity grew 1.33 per cent a year in 1954--63 and 0.53 per cent a year in 2013--23, with no upward trend in between. Meanwhile research effort
rose: researchers per million people grew by half in the United States between 1996 and 2022, doubled in Germany, and more
than quadrupled in Korea and China. A model in which growth rises with the number of researchers must explain why the right
panel rose while the left panel did not.

\begin{recall}
From the school volumes: ideas, institutions and growth, and creative destruction, the replacement of old products and
firms by new ones. From the Year 1 volume: the Solow model and total factor productivity. From \cref{ch:b2:externalities-and-public-goods-theory}: rivalry and
excludability. From \cref{ch:b2:ramsey-model}: balanced growth and the steady state.
\end{recall}

\begin{notation}[Symbols of this chapter]\label{not:b2:endogenous-growth:symbols}
$A$ is the stock of ideas (technology); $L$ the workforce, growing at rate $n$; $L_A = s_RL$ the researchers, $s_R$ the research share; $\nu$ the productivity of research;
$\phi$ the strength of the effect of existing ideas on new ones; $g_x = \dot x/x$ the growth rate of $x$.
\end{notation}

\section{Constant returns to capital}\label{sec:b2:endogenous-growth:ak}

\begin{definition}[AK model]\label{def:b2:endogenous-growth:ak}
The \emph{AK model}\index{AK model} is a growth model in which output is proportional to a broad measure of capital, $Y = AK$,
so that capital faces no diminishing returns and its growth does not slow down.
\end{definition}

With saving rate $s$ and depreciation $\delta$, $\dot K = sAK - \delta K$, so $g_Y = g_K = sA - \delta$: growth rises with saving for ever. With $A = 0.5$ and $\delta = 0.05$, a saving rate of 20 per cent gives growth of 5 per cent a
year and a saving rate of 25 per cent gives 7.5 per cent.

\begin{omfigure}
\begin{tikzpicture}[econ]
\begin{axis}[omrecord, grid=none, xmin=0, xmax=0.4, ymin=-0.06, ymax=0.16, width=0.5\linewidth, height=0.4\linewidth,
    xlabel={Saving rate, $s$}, ylabel={Growth rate}, scaled x ticks=false, scaled y ticks=false,
    xticklabel style={/pgf/number format/fixed}, yticklabel style={/pgf/number format/fixed}]
  \addplot[omDef, very thick, domain=0:0.4, samples=2] {0.5*x - 0.05};
  \draw[black!40] (axis cs:0,0) -- (axis cs:0.4,0);
  \node[eqpoint] at (axis cs:0.2,0.05) {};
  \node[eqpoint] at (axis cs:0.25,0.075) {};
  \draw[guide] (axis cs:0.2,0.05) -- (axis cs:0.2,-0.06);
  \draw[guide] (axis cs:0.25,0.075) -- (axis cs:0.25,-0.06);
\end{axis}
\end{tikzpicture}

{\small Growth in the \omterm{def:b2:endogenous-growth:ak}{AK model}, $g = sA - \delta$, with $A = 0.5$ and $\delta = 0.05$: each extra point of saving adds half a point of growth, permanently.}
\end{omfigure}

\begin{model}{AK model}
\assumes Constant returns to a broad capital stock that includes knowledge and skills, so that accumulation never runs
into diminishing returns.
\predicts Permanent differences in growth rates between economies that save at different rates; a permanent change in
the saving rate changes growth permanently.
\tested On American time series: persistent changes in the determinants the model names should appear as persistent
changes in the growth rate (Jones 1995).
\breaks American growth rates show no large persistent changes, and many AK-style models fail that test.
\end{model}

\section{Ideas used by all at once}\label{sec:b2:endogenous-growth:nonrival}

\begin{definition}[Nonrival good]\label{def:b2:endogenous-growth:nonrival}
A \emph{nonrival good}\index{nonrival good} is a good whose use by one person does not reduce the amount available to
others.
\end{definition}

A design, a formula or a method is nonrival: once found, it can be used by every worker in every firm at once. This has
two consequences. Production with ideas has increasing returns: doubling workers and machines doubles output with the same
ideas, so doubling ideas as well more than doubles it. And a market in which ideas are sold at marginal cost, which is near
zero, cannot pay for their discovery. Some excludability, through patents or secrecy and the market power they give, is
needed for private research to pay; the models of this chapter therefore have monopolistic firms selling designs at a
price above marginal cost.

\section{The research-sector model}\label{sec:b2:endogenous-growth:research}

\begin{definition}[Research sector]\label{def:b2:endogenous-growth:research-sector}
The \emph{research sector}\index{research sector} of a growth model is the part of the economy that employs workers to
produce new ideas rather than goods.
\end{definition}

\begin{definition}[Ideas production function]\label{def:b2:endogenous-growth:ideas-production}
An \emph{ideas production function}\index{ideas production function} gives the flow of new ideas as a function of the
researchers employed and the existing stock of ideas: $\dot A = \nu L_AA^\phi$.
\end{definition}

Goods are produced with ideas and the remaining workers: $Y = A(1 - s_R)L$, so output per worker is $y = A(1 - s_R)$ and its growth rate equals that of $A$. With $\phi = 1$,
standing on the shoulders of giants exactly offsets the fishing out of easy ideas, and $g_A = \dot A/A = \nu s_RL$: growth is proportional to the number of
researchers. A society that doubles its research share doubles its growth rate, for ever.

\begin{definition}[Scale effect]\label{def:b2:endogenous-growth:scale-effect}
A growth model has a \emph{scale effect}\index{scale effect (growth models)} when its long-run growth rate rises with the
size of the economy, usually through the number of researchers.
\end{definition}

\begin{model}{The research-sector model}
\assumes Ideas produced by researchers with $\phi = 1$, nonrival ideas sold by firms with market power, a constant share of
workers in research.
\predicts Long-run growth proportional to the number of researchers; larger economies grow faster.
\tested On the post-war record of research effort and growth in the United States (Jones 1995; the opening chart).
\breaks Researchers multiplied while American growth stayed flat: the time series reject these models particularly
strongly.
\end{model}

\section{The scale effect and its refutation}\label{sec:b2:endogenous-growth:scale}

\begin{record}[The scale effect and the post-war data]
% ledger: b2:jones1995
Jones (1995) put the endogenous growth models to a simple test. In these models, permanent changes in variables such as
the number of researchers or the investment rate change the growth rate permanently. American growth rates show no large
persistent changes. So the variables the models name must show no persistent changes either, or their changes must
offset. Many AK-style models and models with a \omterm{def:b2:endogenous-growth:research-sector}{research sector} fail this test, the latter particularly
strongly. The opening chart shows the same pattern in recent decades: research effort up, productivity
growth not up.
\end{record}

\section{The semi-endogenous version}\label{sec:b2:endogenous-growth:semi}

\begin{definition}[Semi-endogenous growth model]\label{def:b2:endogenous-growth:semi-endogenous}
A \emph{semi-endogenous growth model}\index{semi-endogenous growth model} is a growth model with a \omterm{def:b2:endogenous-growth:research-sector}{research sector} in
which ideas become harder to find as they accumulate, $\phi < 1$, so that long-run growth depends on the growth of the number of
researchers, not on its level or share.
\end{definition}

With $\phi < 1$ and a constant research share, $g_A = \nu s_RLA^{\phi - 1}$. Along a balanced growth path $g_A$ is constant, so $L$ and $A^{1-\phi}$ grow at the same rate: $n = (1 - \phi)g_A$, and
\begin{equation}\label{eq:b2:endogenous-growth:semi}
g_A = \frac{n}{1 - \phi}.
\end{equation}
With $n = 1$ per cent and $\phi = 0.5$, ideas and output per worker grow 2 per cent a year. Policy that raises the research share does not change this rate:
it raises growth for a time and leaves a higher level of output per worker.

\begin{omfigure}
\begin{tikzpicture}[econ]
\begin{axis}[omrecord, grid=none, xmin=0, xmax=10, ymin=0, ymax=0.06, width=0.55\linewidth, height=0.4\linewidth,
    xlabel={Stock of ideas, $A$ (researchers held fixed)}, ylabel={$g_A$}, scaled y ticks=false,
    yticklabel style={/pgf/number format/fixed}, legend style={at={(1.03,0.5)}, anchor=west}]
  \addplot[omDef, very thick, domain=0:10, samples=2] {0.03};
  \addplot[omThm, very thick, domain=0.25:10, samples=80] {0.03/sqrt(x)};
  \legend{$\phi = 1$, $\phi = 0.5$}
\end{axis}
\end{tikzpicture}

{\small The growth of ideas with a fixed number of researchers. With $\phi = 1$ it does not fall as ideas accumulate; with $\phi < 1$ it falls, and sustained
growth needs a growing number of researchers.}
\end{omfigure}

The listing doubles the research share, from 2 to 4 per cent, in an economy on its balanced growth path, and follows growth
and output per worker relative to the economy that kept 2 per cent.

\omcode{code/bachelor-2/24-endogenous-growth/listings/ideas.py}{1}{28}{Doubling the research share in the
semi-endogenous model.}\label{lst:b2:endogenous-growth:ideas}

\omoutput{code/bachelor-2/24-endogenous-growth/outputs/ideas.txt}{Growth jumps to 4 per cent and falls back toward 2;
output per worker ends almost four times higher, but only after centuries.}\label{out:b2:endogenous-growth:ideas}

% code: listings/ideas.py
At first output per worker falls by 2 per cent, since workers leave goods production for research. Growth doubles, then falls back: after
50 years it is still 2.87 per cent, and the level of output per worker is 1.9 times the baseline. In the long run the level gain is $2^{1/(1-\phi)} \times 0.96/0.98 = 3.918$, reached only
over centuries. The semi-endogenous model fits flat growth with rising research effort; its price is that long-run growth
is tied to population growth, which policy does not usually aim at.

\begin{model}{Semi-endogenous growth}
\assumes A \omterm{def:b2:endogenous-growth:research-sector}{research sector} in which ideas become harder to find ($\phi < 1$), and a workforce that grows.
\predicts Long-run growth $n/(1 - \phi)$ independent of policy; research subsidies raise levels, not long-run growth; \omterm{def:b2:endogenous-growth:research-productivity}{research productivity}
falling while research effort rises.
\tested Consistent with flat American growth alongside rising research effort, and with falling \omterm{def:b2:endogenous-growth:research-productivity}{research productivity}
(Bloom, Jones, Van Reenen and Webb 2020).
\breaks Long-run growth tied to population growth implies that growth ends when population stops growing, a prediction
the data cannot yet test.
\end{model}

\section{Growth through replacement}\label{sec:b2:endogenous-growth:replacement}

Another family of models makes growth come from new products that replace old ones: each innovation makes the previous
version obsolete, and the innovator earns the profits of the market until the next innovation replaces it. This is
the creative destruction of the school volumes, in formal dress. Two features distinguish these models. Growth destroys
firms and jobs as well as creating them, so the gains from growth come with reallocation. And an innovator does not count
the profits it takes from the firm it displaces, so private research may be larger or smaller than the level that
maximises welfare: the business-stealing effect pushes it up, the fact that innovators capture only part of the value of
the ideas later research builds on pushes it down.

\section{Are ideas getting harder to find?}\label{sec:b2:endogenous-growth:harder}

\begin{definition}[Research productivity]\label{def:b2:endogenous-growth:research-productivity}
\emph{Research productivity}\index{research productivity} is the growth rate of ideas divided by the number of
researchers producing them, $g_A/L_A$.
\end{definition}

\begin{record}[Ideas are getting harder to find]
% ledger: b2:bloom2020
Bloom, Jones, Van Reenen and Webb (2020) measured \omterm{def:b2:endogenous-growth:research-productivity}{research productivity} directly, field by field. Their best-known case is
computer chips: keeping chip density doubling every two years required more than 18 times as many researchers in their latest
years as in the early 1970s. Across their cases, research effort rose substantially and \omterm{def:b2:endogenous-growth:research-productivity}{research productivity} fell
sharply: everywhere they looked, ideas were getting harder to find. Growth was sustained by putting ever more
people to work on it, the pattern the semi-endogenous model predicts.
\end{record}

If growth is held constant while researchers multiply by 18, \omterm{def:b2:endogenous-growth:research-productivity}{research productivity} falls to one eighteenth. If growth halves while researchers
quadruple, it falls to an eighth. In the semi-endogenous model this is the expected state of affairs, not a crisis: $\phi < 1$ means each idea makes
the next one harder.

\section{Exercises}\label{sec:b2:endogenous-growth:exercises}

\begin{exercise}[$\star$]\label{exo:b2:endogenous-growth:1}
Read the opening chart. In which decade did American productivity grow fastest, and in which slowest? By what factor did
researchers per million people rise in each country?
\end{exercise}

\begin{exercise}[$\star$]\label{exo:b2:endogenous-growth:2}
In the \omterm{def:b2:endogenous-growth:ak}{AK model} with $A = 0.5$ and $\delta = 0.05$, compute growth at saving rates of 20 and 25 per cent.
\end{exercise}

\begin{exercise}[$\star$]\label{exo:b2:endogenous-growth:3}
Classify as rival or nonrival: a recipe, a meal, a computer program, a computer, a vaccine's formula, a dose of vaccine.
\end{exercise}

\begin{exercise}[$\star$]\label{exo:b2:endogenous-growth:4}
State an observation that would falsify the research-sector model with $\phi = 1$. Has it been made?
\end{exercise}

\begin{exercise}[$\star\star$]\label{exo:b2:endogenous-growth:5}
Predict, then compute: with $n = 1$ per cent, what is long-run growth when $\phi = 0.75$? Compare with $\phi = 0.5$.
\end{exercise}

\begin{exercise}[$\star\star$]\label{exo:b2:endogenous-growth:6}
With $\phi = 0.5$, the research share rises from 2 to 3 per cent. Find the long-run ratio of output per worker to the baseline.
\end{exercise}

\begin{exercise}[$\star\star$]\label{exo:b2:endogenous-growth:7}
Growth of ideas halves while the number of researchers quadruples. By what factor does \omterm{def:b2:endogenous-growth:research-productivity}{research productivity} change?
\end{exercise}

\begin{exercise}[$\star\star$]\label{exo:b2:endogenous-growth:8}
Modify the listing to set $\phi = 0.75$. What changes in the long-run growth rate, the level gain, and the speed of the transition?
\end{exercise}

\begin{exercise}[$\star\star$]\label{exo:b2:endogenous-growth:9}
Explain why a market in which designs sell at marginal cost cannot pay for research, and why patents trade one inefficiency
for another.
\end{exercise}

\begin{exercise}[$\star\star\star$]\label{exo:b2:endogenous-growth:10}
Derive \eqref{eq:b2:endogenous-growth:semi} from $\dot A = \nu s_RLA^\phi$ with a constant $s_R$, assuming growth of $A$ is constant.
\end{exercise}

\begin{exercise}[$\star\star\star$]\label{exo:b2:endogenous-growth:11}
Let $z = g_A$. Show that $\dot z = z(n - (1 - \phi)z)$, and solve this equation from $z_0$. With $n = 0.01$ and $\phi = 0.5$, how long does growth take to fall from 4 to 3 per cent?
\end{exercise}

\begin{exercise}[$\star\star\star$]\label{exo:b2:endogenous-growth:12}
On a balanced growth path, show that $A^{1-\phi} = \nu s_RL/g_A$, and hence that doubling $s_R$ multiplies the long-run level of $A$ by $2^{1/(1-\phi)}$.
\end{exercise}

\section{Problem: Brenland's research labs}\label{sec:b2:endogenous-growth:problem}

\begin{problem}[{Weekend problem --- Brenland's research labs: the ideas function, balanced growth, doubling the research
share, and the gain in the level of income}]
\label{pb:b2:endogenous-growth:1}
Brenland's ideas grow as $\dot A = \nu s_RLA^{0.5}$; its workforce grows 1 per cent a year; 2 per cent of workers do research; output per worker is $A(1 - s_R)$.

\medskip
\noindent\textbf{Part I --- The ideas function.}
\begin{enumerate}
  \item What is $\phi$, and what does $\phi < 1$ say about research?
  \item Write the growth rate of ideas.
  \item With researchers held fixed, what happens to $g_A$ as $A$ grows?
  \item Why must the number of researchers grow for growth to continue?
  \item Define \omterm{def:b2:endogenous-growth:research-productivity}{research productivity} in this model.
\end{enumerate}

\medskip
\noindent\textbf{Part II --- Balanced growth.}
\begin{enumerate}[resume]
  \item What is the long-run growth rate of ideas?
  \item Of output per worker?
  \item Would it differ if 5 per cent of workers did research?
  \item What would it be if the workforce stopped growing?
  \item What would it be with $\phi = 0.75$?
  \item If $\phi$ were 1, what would happen to growth as the workforce grows?
\end{enumerate}

\medskip
\noindent\textbf{Part III --- Doubling the research share.} Brenland raises $s_R$ to 4 per cent.
\begin{enumerate}[resume]
  \item What happens to output per worker at once?
  \item What happens to the growth rate of ideas at once?
  \item Write the equation for $g_A$ over time, and find how long it takes to fall to 3 per cent.
  \item What is growth after 50 years?
  \item What is the long-run growth rate?
  \item Why does growth fall back?
\end{enumerate}

\medskip
\noindent\textbf{Part IV --- The level effect.}
\begin{enumerate}[resume]
  \item By what factor does the long-run level of $A$ rise?
  \item And output per worker, after the cost of the extra researchers?
  \item How much of the gain is reached after 50 years (use the listing)?
  \item With $\phi = 0.75$, what would the long-run level gain be?
  \item Why might Brenland choose to double research even though long-run growth is unchanged?
  \item What observation would show that Brenland's $\phi$ is larger than 0.5?
  \item State the result: the balanced growth rate and the long-run level gain from doubling the research share.
\end{enumerate}
\end{problem}
